
\section{Finite relations, backward viability, and stable erasure}
\label{sec:relational}

\subsection{Why relations are first-class}

Constructive formalization often fails at the point where a proof informally
says ``choose one.''  The issue is not finite search itself: bounded least
search can choose the least code satisfying a decidable predicate.  The issue
is that a witness valid at one precision may not refine to the least witness
at a finer precision, and an exact global branch may not exist at collisions.

PCFM therefore distinguishes:

\begin{itemize}
\item \emph{local selection}: choose the least member of one finite nonempty
set;
\item \emph{finite coherent selection}: choose a path through a fixed finite
sequence of compatibility relations;
\item \emph{infinite coherent selection}: choose a path through every
precision level.
\end{itemize}

The first two are primitive-recursive when explicit bounds are given.  The
third is a genuinely stronger compactness or choice principle and is not
silently inferred.

\subsection{Backward viability}

Let time range over \(0,\ldots,T\).  At time \(t\), candidates are codes
\(a<B_t\).  Let
\[
  \Good_t(a)\in\Bool,\qquad
  R_t(a,b)\in\Bool
\]
be reflected predicates, with \(R_t\) relating a candidate at \(t\) to one at
\(t+1\).  Let \(\mathsf{Terminal}(a)\) be the final acceptance predicate.

Define viability backward:
\[
\begin{aligned}
 \Viable_T(a)
   &\defeq \Good_T(a)\wedge\mathsf{Terminal}(a),\\
 \Viable_t(a)
   &\defeq \Good_t(a)\wedge
      \exists b<B_{t+1}\,
        (R_t(a,b)\wedge\Viable_{t+1}(b)).
\end{aligned}
\]
Every existential is compiled by bounded search, so each \(\Viable_t\) has a
Boolean characteristic term.  The complete table is computed by a bounded
reverse fold.

\begin{theorem}[Finite path extraction]
If \(\exists a<B_0\,\Viable_0(a)\), then a primitive-recursive selector returns
a path \(a_0,\ldots,a_T\) such that
\[
  \Good_t(a_t)=1,\qquad
  R_t(a_t,a_{t+1})=1,\qquad
  \mathsf{Terminal}(a_T)=1.
\]
iPRA proves each property of the selected path.
\end{theorem}

\begin{proof}
Choose \(a_0\) as the least viable initial candidate.  Given a selected viable
\(a_t\), the definition of \(\Viable_t(a_t)\) includes a bounded existential
for a viable compatible successor.  Apply bounded least witness to obtain
\(a_{t+1}\).  Primitive recursion on \(t<T\) constructs the finite path.

For the object proof, the invariant says that the path prefix has valid nodes
and edges and that its last node is viable.  This is a bounded Boolean
equation.  The base follows from the initial search.  At the step, the
least-witness specification supplies the next edge and viability.  At time
\(T\), final viability yields terminal acceptance.
\end{proof}

\begin{corollary}[Finite dependent choice]
For every fixed primitive-recursive horizon \(T\), bounded totality of each
relation on viable nodes implies the existence of a compatible path of length
\(T\).
\end{corollary}

This is the exact finite principle needed by theorem-specific computations.
It is not countable dependent choice.

\subsection{Failure certificates}

If no viable initial state exists, the backward table is itself a finite
refutation certificate.  At a failed node it records either:
\[
   \Good_t(a)=0
\]
or that every bounded successor is incompatible or nonviable.  Boolean
reflection turns this into an object proof of failure.

Thus a bounded game or search can return dual evidence:
\begin{itemize}
\item a successful path;
\item a backward table proving that no successful path exists.
\end{itemize}
This is particularly useful for exact finite classifications and
computer-generated counterexample searches.

\subsection{Exact witness erasure}

\begin{theorem}[Exact stable erasure]
Suppose \(W(z)\) is a bounded nonempty witness predicate and \(O(z)\) a
primitive-recursive observation.  If iPRA proves
\[
  W(z)=1\wedge W(z')=1\imp O(z)=O(z'),
\]
then the least witness
\[
  z_*=\mu z<B\,[W(z)=1]
\]
satisfies \(W(z_*)=1\), and \(O(z_*)\) is independent of the selection
algorithm.  Any theorem proved from \(O(z_*)\) may be stated as a theorem
about the witness relation rather than the chosen representative.
\end{theorem}

\begin{proof}
Bounded least witness gives \(W(z_*)=1\).  For any other valid \(z\), apply the
pairwise stability premise to \(z_*,z\).  Equality substitution transports
every conclusion depending only on \(O\).
\end{proof}

This theorem handles symmetric polynomial data of a root cloud, a determinant
computed from different row-reduction traces, a quotient class represented by
different codes, and any finite certificate whose observable output is
canonical.

\subsection{Quotient-stable erasure}

Let the output presentation \(C\) have equivalence
\(\Eqv_C\).  Replace equality in the preceding theorem by
\(\Eqv_C(O(z),O(z'))=1\).  Every downstream map must be extensional for
\(\Eqv_C\).  Composition then yields a well-defined quotient-level result
without forming a quotient object in iPRA.

This is the recommended treatment of:
\begin{itemize}
\item rational numbers before denominator normalization;
\item finite sets represented by lists;
\item localizations represented by numerator--denominator pairs;
\item modules or ideals represented by finite generating lists;
\item homology classes represented by cycles modulo boundaries.
\end{itemize}

\subsection{Error-stable erasure}

Suppose outputs are rational or integer-coded metric data and iPRA proves
\[
  W(z)=W(z')=1\imp d(O(z),O(z'))\le\varepsilon.
\]

\begin{theorem}[Diameter erasure]
The least selected output \(O(z_*)\) is within \(\varepsilon\) of every valid
output.  If a downstream observer \(F\) has a certified modulus
\[
  d(x,y)\le\varepsilon\imp
  d'(F(x),F(y))\le\delta,
\]
then \(F(O(z_*))\) is \(\delta\)-stable across all witnesses.
\end{theorem}

\begin{proof}
The first assertion is pairwise stability with \(z_*\) as one endpoint.  The
second follows by specializing the modulus and composing inequalities.
All rational inequalities are reduced to integer cross multiplication.
\end{proof}

Repeated composition accumulates errors according to explicit algebraic
bounds.  The compiler should never use the phrase ``small enough'' without
computing the requested precision backward.

\subsection{Branch collapse under a gap}

Let each valid witness have a label \(\ell(z)\) in a finite discrete set.
Suppose distinct labels imply metric separation at least \(g>0\), while every
two valid outputs have distance \(<g\).

\begin{lemma}[Gap-based branch collapse]
All valid witnesses have the same label.
\end{lemma}

\begin{proof}
If two labels differed, the separation theorem would give distance at least
\(g\), contradicting the strict diameter bound.  Equality of finite labels is
decidable, so the argument is intuitionistic.
\end{proof}

This is the correct way to recover a named branch of roots, eigenvalues,
connected components, signs, or orientations.  The gap is an explicit guard
or producer obligation.  At a collision the guard fails, and only a
symmetric or relational result is exported.

\subsection{Why horizonwise viability does not give an infinite selector}

\begin{barrier}[No automatic coherent limit]
From
\[
  \forall T\ \exists (a_0,\ldots,a_T)\ \text{compatible}
\]
one may not infer a primitive-recursive infinite sequence
\((a_t)_{t\in\Nat}\) that is compatible at every level.
\end{barrier}

The least path for horizon \(T+1\) can disagree with the least path for horizon
\(T\) at early stages.  Producing a coherent infinite path is a form of
compactness or choice.  PCFM offers three legitimate alternatives:

\begin{enumerate}
\item keep the theorem finite-horizon;
\item prove a stabilization or unique-extension modulus;
\item prove that the requested observable is stable, so no coherent path is
needed.
\end{enumerate}

\subsection{Root clouds as the canonical example}

For a monic polynomial with approximate coefficients, individual roots may
collide and exchange labels.  Represent the answer at precision \(k\) by a
finite multiset of rational boxes covering the root multiset, with certified
Hausdorff or matching error.  Symmetric polynomials, root counts in regions
separated from the boundary, and resultants are stable observers.  A named
root branch requires an isolation disc or a separation gap.

This pattern resolves many difficulties more honestly than forcing an exact
global enumeration of roots.  It also fits the checked finite algebra
infrastructure: the certificate consists of integer polynomial identities,
winding or sign data on a finite boundary, and bounded inequalities.
